\chapter{Introduction}


The prediction of financial processes and the modeling of their 
distribution's properties is
one of the central challenges in quantitative finance.
Traditional forecasting methods generally focus on estimating a single future value,
such as the expected return or the most likely price at a given forecasting horizon.
While point forecasts are useful in many contexts,
they provide only limited information regarding the uncertainty
associated with future market movements.
In practice, many financial applications, including derivative pricing, portfolio allocation, Value-at-Risk estimation and stress testing,
depend not only on the expected future price but on the complete probability distribution of future outcomes.\\

For this reason, distributional forecasting has attracted increasing attention in both financial econometrics and machine learning.
Rather than predicting a single value, 
probabilistic models estimate the entire conditional distribution of a given process given the available historical information.
Such models provide richer information about the likelihood of extreme events, asymmetry, tail risk, and forecast uncertainty,
thereby supporting more informed decision-making under uncertainty.\\

The rapid development of deep learning has introduced a new 
family of probabilistic models capable of approximating highly complex data distributions,
among these, the Generative Adversarial Networks (GANs) have emerged as one of the most successful generative modeling techniques.
Originally developed for image synthesis, GANs learn an implicit probability distribution through the adversarial interaction between two neural networks:
a generator, responsible for producing synthetic samples, and a discriminator (or critic),
responsible for distinguishing generated samples from real observations.
Unlike traditional likelihood-based methods, GANs do not require an explicit specification of the underlying probability density,
making them particularly attractive for modeling complex stochastic processes.\\


Recent research has extended GANs to time-series generation and forecasting through conditional architectures,
where the generator produces future observations conditioned on historical data.\\

The objective of this thesis is to investigate the use of generative models for distributional forecasting
of financial processes and, 
in particular, to design a Conditional Generative Adversarial Network (CGAN) which is able to generate 
the future probability distribution of the price of an asset given its past trajectory, enabling a wide range of financial applications that cannot be addressed through point forecasts alone.
Specifically, distributional forecasts provide the probabilistic information required for derivative pricing, market risk measurement, stress testing, and portfolio optimization.\\

In quantitative risk management, the predicted distribution of future asset prices 
can be directly employed to estimate risk measures such as Value-at-Risk (VaR) and Expected Shortfall (ES), 
allowing financial institutions to quantify the probability and magnitude of adverse market movements without relying on restrictive distributional assumptions.
Similarly, probabilistic forecasts provide a natural framework for scenario generation and stress testing, 
enabling the simulation of realistic future market conditions for regulatory compliance and internal capital assessment.
Another important application is related to derivative pricing, in fact, under the risk-neutral valuation framework, 
the price of a financial derivative is given by the expected discounted value of its future payoff under the probability distribution of the underlying asset. 
Consequently, an accurate estimate of the conditional distribution of future asset prices provides the information required to price European and, more generally, path-dependent derivatives through Monte Carlo simulation.\\

The work begins by discussing simple Conditional GAN architectures capable of directly generating future probability distributions conditioned on historical asset trajectories.
Alternative formulations based on the Wasserstein distance are then explored to improve training stability and mitigate well-known issues such as mode collapse and vanishing gradients. 
Subsequently, recurrent forecasting architectures based on the ForGAN framework are investigated, where probability distributions are recovered through repeated stochastic sampling of future values.
Finally, a novel Scoring Rule ForGAN architecture is proposed, which exploits proper probabilistic scoring rules to improve calibration and distributional accuracy.\\

The proposed methodologies are evaluated using synthetic datasets generated from stochastic financial models. 
Initially, experiments are conducted assuming assets' dynamics can be modelled under the Black and Scholes model assumptions by using Geometric Brownian Motion, 
where analytical probability distributions are available and provide an exact benchmark for model evaluation.
To assess the robustness of the proposed approach under more realistic market dynamics, the experimental framework is subsequently extended by implementing a Heston stochastic volatility simulator, 
allowing the models to be tested on processes exhibiting time-varying volatility.\\

Model performance is assessed through several statistical measures that compare the generated and theoretical probability distributions, including Jensen-Shannon divergence, Hellinger distance, Total Variation distance, Wasserstein distance, and the Kolmogorov-Smirnov goodness-of-fit test.
These metrics provide complementary perspectives on distributional accuracy, calibration, and the ability of the proposed models to reproduce the statistical properties of financial processes.

The remainder of this thesis is organized as follows: chapter 2 is delegated to an analysis of the literature about generative models and in particular the Generative Adversarial Networks. It introduces the Conditional GAN and
Conditional Wasserstein GAN architectures and their mathematical formulation, highlighting the main challenges and limitations of these models in the context of distributional forecasting in the financial domain; 
chapter 3 investigates the ForGAN framework and its Wasserstein variant for probabilistic forecasting; 
chapter 4 proposes the Scoring Rule ForGAN architecture and analyzes its performance relative to previous models;
chapter 5 introduces the Heston stochastic volatility simulator and evaluates the proposed methodology under a more realistic financial setting. 
Finally, the concluding chapter summarizes the main findings, discusses the limitations of the proposed approach, and outlines directions for future research.
